\chapter{The Euler--Lagrange Equivalence}
\label{chap:euler-lagrange-equivalence}

The previous chapter brought a path to the edge of stationarity: a path is stationary when its first
variation vanishes, and flatness is the clean, exact form of that vanishing. This chapter takes the
last step, from stationarity to the equation that physics writes for it --- the Euler--Lagrange
equation --- and in taking it the construction arrives at the one place in the entire book where it
cannot proceed by building. Everything until now has been exhibited: every number constructed, every
theorem proved by producing a witness, every characteristic inhabited. Here, for the first and only
time, the construction reaches a question it cannot answer by computing, and instead of pretending
otherwise it does the honest thing --- it names the assumption, grants it openly, and marks exactly
where the building stops. The long constructive discipline of Part I exists so that this single
assumption, when it finally appears, has nowhere to hide.

\section{The equivalence, proved}
\label{se:the-equivalence-proved}

Start with what is solid, because most of the chapter is. The construction forms the discrete
derivative of its action --- the finite, exact analogue of differentiating a functional --- and writes
down the discrete Euler--Lagrange equation: the statement that this derivative's residual vanishes in
every variation direction. Then it proves the equivalence the chapter is named for. A path is
stationary if and only if it satisfies the discrete Euler--Lagrange equation. Both directions are
genuine theorems. Stationarity implies the equation, because the action's change in any direction
decomposes into a part the derivative captures and a remainder the construction shows is identically
zero, so flatness forces the derivative to vanish. And the equation implies stationarity, by running
the same decomposition in reverse. Neither direction assumes anything; each is established by
exhibiting the algebra that connects the two descriptions.

It is worth seeing the shape of that proof, because it explains why this direction costs nothing while
the next costs everything. The change in the action along a variation splits cleanly into two pieces: a
part the discrete derivative is defined to capture --- the leading response, the analogue of a
derivative times a displacement --- and a remainder, everything the leading part misses. The whole
point of the construction's setup is that, for its finite gauge action, this remainder is identically
zero: the discrete derivative captures the change exactly, with nothing left over. Once that is in
hand, the equivalence is immediate in both directions. If the path is flat, the total change is zero,
and since the remainder is already zero the derivative's part must be zero too --- the Euler--Lagrange
equation. If the derivative's part is zero, then since the remainder is zero the total change is zero
in that direction --- flatness. The equivalence is exact precisely because the remainder is exactly
zero, and that exactness is what a finite construction can deliver. What it cannot deliver comes next.

This equivalence is the formal heart of the variational method, and it is worth seeing what it buys.
Stationarity is a global, geometric condition --- the path sits where the action is flat. The
Euler--Lagrange equation is a local, algebraic one --- a derivative vanishes. The theorem says these
are the same condition wearing different clothes, so a problem stated as ``find the flat path'' can be
solved as ``solve this equation,'' which is the move that makes the calculus of variations a
calculational tool rather than a description. The construction tanges the solutions by the trait the
equation tests --- the residual vanishes here --- and the equivalence guarantees that the paths so
selected are exactly the stationary ones. The experiment under the name of minimizing variations is
the classical original of this: form the functional whose value is an integral over a path, take its
first variation under a fixed-endpoint perturbation, set it to zero, and the Euler--Lagrange equation
falls out as the condition the integrand must satisfy. The construction's finite version reaches the
same equation by the same logic, with exact differences in place of derivatives.

A concrete reading anchors the abstraction. A free particle, in the classical theory, has an action
that totals its kinetic energy along a path, and the Euler--Lagrange equation its first variation
yields is simply that its acceleration is zero --- it moves in a straight line at constant speed. The
global statement, that the path is the one making the action stationary, and the local statement, that
the acceleration vanishes at every instant, are by the equivalence the same statement: one checked by
comparing whole paths, the other by solving an equation point by point. The construction's discrete
version replaces the instants with a finite mesh and the acceleration with a discrete second
difference, but the equivalence is the same, and it is this that lets the next chapter solve real
problems --- it never has to search the space of paths, only to solve the equation the equivalence
hands it.

\section{The one thing it cannot decide}
\label{se:the-one-thing-it-cannot-decide}

Now the catch, and it is the most important paragraph in the book. The equivalence is proved, but to
\emph{use} it --- to decide whether a given path actually is stationary --- you must check that its
first variation vanishes. Not in one direction, not in a thousand: in \emph{every} variation
direction. And the space of variation directions is infinite. The construction can compute the change
of the action in any single direction it is handed, exactly and finitely; what it cannot do is range
over all directions at once and certify that the change is zero in every one of them. Stationarity is
a universal claim quantified over an infinite domain, and a universal over an infinite domain is not a
computation. No finite procedure can check infinitely many cases, and the construction has spent ten
chapters refusing to pretend otherwise.

It is worth being exact about the asymmetry, because it is the crux. Testing one variation is a finite
act: hand the construction a direction, and it computes the action's change along it and reports
whether that change is zero, all in bounded time. Testing every variation is not a finite act, and no
amount of cleverness makes it one, because the variations are not a list that can be walked to its end
--- they form an infinite space, and a property that must hold throughout an infinite space is a
universally quantified claim, not a value to be computed. This is the precise boundary of the whole
construction. Everything it has built has been, at bottom, a finite procedure returning a verdict;
stationarity is the first object it needs that is not of that kind. It is a claim \emph{about} all the
finite procedures at once, and that is a different and larger thing than any one of them.

So here it stops building, and grants. At precisely the point where deciding stationarity would require
the one thing a finite construction cannot supply --- a completed survey of infinitely many directions
--- the construction introduces a single axiom. It consults an oracle: it assumes, without proof and
without apology, that the question ``is this path stationary?'' has a decidable answer, even though it
cannot compute that answer itself. This one axiom is the sole assumption of its kind in the entire
development. Everything else was built; this alone is granted. And the construction is precise about
what the assumption \emph{is} in physical terms. The decidability of stationarity is the law of
motion: the claim that nature, faced with a path, does settle whether that path is the one it takes.
The oracle is the construction's way of saying that this settling happens --- that there is a fact of
the matter about which path is stationary --- while admitting in the same breath that establishing the
fact is the one query nature answers and a finite procedure cannot.

What the oracle supplies is narrow and specific, and saying exactly what it is keeps the assumption
from sprawling. It does not assert that the path is stationary; it asserts only that the question has a
definite answer --- that stationarity, for any given path, is either true or false and can be treated
as decided. The construction still computes everything it can: it forms the action, the derivative, the
direct residual, the equivalence. The oracle enters at exactly one point, supplying the single bit the
construction cannot produce --- whether the residual vanishes across all directions --- and from that
one granted bit the rest follows by ordinary computation. It is the smallest assumption that does the
job: not a path handed down, not a law postulated in full, but one decision, about one universal,
granted because the alternative is to pretend a finite survey can be infinite.

\begin{quote}
\emph{A note on the labels --- and the one assumption.} This is the chapter to read slowly. The
construction has, to here, assumed almost nothing: proposition extensionality, the quotient that
collapses observationally equal readings, and the bridges of physical naming, all declared as they
arose. The Euler--Lagrange oracle is different in kind. It is not a naming and not a convenience; it is
a genuine logical assumption, the assertion that a universal quantifier over an infinite variation
space can be decided when no finite computation can decide it. It is the one
non-constructive grant the whole development needs --- the assertion that this universal, which no
finite procedure can survey, is nonetheless \emph{decided} --- and it wears, in physics, the name of the
law of motion: the single query that nature answers and the finite construction cannot. The entire constructive discipline of the first part ---
doing without the excluded middle, without a free power to choose, building every other result by
exhibiting a witness --- exists to make this one assumption stand out in full relief. It is the load
the whole structure hangs from, and the honest thing is to point at it directly: \emph{here, and only
here, the construction stops proving and starts trusting.}
\end{quote}

It helps to see why the wall is exactly where it is, because it is the same wall the book has hit
before. To confirm a path stationary, the construction can test direction after direction and find
the action flat in each --- the first, the second, the thousandth --- and every test that passes is a
small confirmation, exactly as a study accumulates agreeing trials. But no finite number of passing
tests is the universal claim; the next direction could always be the one where the action is not flat,
and the construction can no more rule that out by testing than Hume's observer could rule out a future
counterexample, or than a finite reader could certify that a stream of bits was lawful rather than
noise. The oracle is the assumption that the survey, could it be completed, would come back clean ---
the leap from ``flat in every direction I have checked'' to ``flat in every direction,'' which is
precisely the leap a finite construction is not entitled to take and which physics takes every time it
writes a law of motion. The construction funges all those directions into a single universal and then
grants, by axiom, that the universal has a definite truth-value. That grant is the whole of what it
assumes here.

With that one grant in place, the rest is bookkeeping the construction does for itself. Given the
oracle's verdict that the residual vanishes across all directions, and its own computation that the
direct residual vanishes too, it forms the total residual and reads off a residue: zero --- the bare
truth that true equals true --- when the path is stationary, and the path's own endpoint, a non-zero
marker, when it is not. The residue is the construction's honest receipt. It records, in one value it
can compare and check, whether the path it was handed obeys the law of motion, and it carries in its
provenance the exact debt it owes --- one consulted oracle, and nothing else. Every later result that
reads a residue of zero is reading, folded into that zero, the single assumption this chapter made and
the long chain of computation that surrounds it.

\section{What is demonstrated, and what is assumed}
\label{se:demonstrated-assumed-ch10}

This chapter's separation of built from assumed is the most important one in the book, so it is worth
stating with no softening.

What is demonstrated is the equivalence itself, and it is fully proved. The discrete Euler--Lagrange
equation is a finite, exact object; the theorem that stationarity holds if and only if the equation
holds is established in both directions by exhibiting the decomposition that connects them; and the
residue the construction reads off a path is zero exactly when the total residual vanishes, which is a
computation it can carry out once the decision is in hand. None of this is assumed. The machinery that
turns ``flat'' into ``solves an equation'' is genuine, checked mathematics.

A word on why the assumption sits here, in the open, rather than dissolved quietly into the machinery.
A less scrupulous construction could have buried the same leap in a dozen places --- assuming, here and
there, that some search terminates or some property holds throughout an infinite range --- and no
single line would have looked like an axiom. This construction does the opposite. It drives every such
need down to one point, the decidability of stationarity, and grants it once, by a named axiom, so that
the whole development's debt is a single declared line a reader can find, weigh, and refuse if she
likes. That localization is itself the achievement of Part I's discipline: not that the construction
assumes little, though it does, but that it assumes its little in one visible place. An honest theory
is not one with no assumptions; it is one whose assumptions you can point to.

What is assumed is a single thing, and it is the keystone: that stationarity --- a universal over the
infinite space of variations --- is decidable, granted by the one custom axiom the development
contains. This is the construction's only appeal to something it cannot compute, and it is deliberately
the law of motion itself, the physical fact that there is a path nature takes. Every honest reckoning
of this book's claims turns on holding that line clearly. The construction demonstrates, with nothing
up its sleeve, that the variational and the differential descriptions of motion are the same; it
assumes, openly and exactly once, that the question those descriptions pose has an answer. The next
chapter turns from this principle to practice --- to actually solving the Euler--Lagrange equation with
splines, with the Lanczos iteration, and over real partitions --- on the strength of the one assumption
named here, and no other.
